\chapter{McNemar 検定}\label{ch:m06}

\begin{priobox}{B}
3回出題（\kakomon{2018}{2}，\kakomon{2019}{3}〔2〕，\kakomon{2024}{1}〔2〕〔3〕）．いずれも手法名は出ないか，ヒント付きである．
しかし「一致ペアは差に寄与しない」「$H_0$ の下の分散は $(n_{12}+n_{21})/n$ で推定する」「不一致ペアを条件とすれば $\Bin(s,1/2)$」「$\chi^2_1$ で判定」を知っていると数分で終わる．
本編の第4章4.6節．
\end{priobox}

% ============================================================
\section{概要}\label{sec:m06-gaiyo}
$n$ 組のペア（または $n$ 人に2つの処置・検査）の2値の結果を次のように集計する．セル度数は多項分布 $\mathrm{Mult}(n;\pi_{11},\pi_{12},\pi_{21},\pi_{22})$ に従う．
\begin{center}\small
\begin{tabular}{@{}lccc@{}}
\toprule
 & 処置2 陽性（有効） & 処置2 陰性（無効） & 計\\
\midrule
処置1 陽性 & $n_{11}$ & $n_{12}$ & $n_{1\cdot}$\\
処置1 陰性 & $n_{21}$ & $n_{22}$ & $n_{2\cdot}$\\
\midrule
計 & $n_{\cdot1}$ & $n_{\cdot2}$ & $n$\\
\bottomrule
\end{tabular}
\end{center}
陽性率の差は $\pi_{1\cdot}-\pi_{\cdot1}=\pi_{12}-\pi_{21}=:\delta$ で，一致ペア（$n_{11},n_{22}$）は差に寄与しない．
\begin{pointbox}[覚える式]
\begin{itemize}
\item McNemar 検定（$H_0:\pi_{12}=\pi_{21}$）：$W=\dfrac{(n_{12}-n_{21})^2}{n_{12}+n_{21}}\ \xrightarrow{d}\ \chi^2_1$．
\item 条件付き二項：不一致ペア数 $s=n_{12}+n_{21}$ を与えると $H_0$ の下で $n_{12}\sim\Bin(s,1/2)$，$z=\dfrac{n_{12}-s/2}{\sqrt{s/4}}$，$z^2=W$．
\item 差の推定：$\hat\delta=\dfrac{n_{12}-n_{21}}{n}$，$\V[\hat\delta]=\dfrac{\pi_{12}+\pi_{21}-(\pi_{12}-\pi_{21})^2}{n}$，
      CI $\hat\delta\pm z_{\alpha/2}\sqrt{\dfrac{n_{12}+n_{21}}{n^2}-\dfrac{(n_{12}-n_{21})^2}{n^3}}$．
\item 尤度比検定：$G^2=2\Bigl[n_{12}\log\dfrac{2n_{12}}{n_{12}+n_{21}}+n_{21}\log\dfrac{2n_{21}}{n_{12}+n_{21}}\Bigr]$，自由度1．
\item 条件付きオッズ比：$n_{12}/n_{21}$．
\end{itemize}
\end{pointbox}

% ============================================================
\section{詳細}\label{sec:m06-shosai}
\subsection{差の平均と分散（\kakomon{2024}{1}の見方）}
被験者 $i$ の処置1・2の結果を $X_i,Y_i\in\{0,1\}$，$\pi_{jk}=\Prob(X_i=j,Y_i=k)$ とし，$Z_i=X_i-Y_i\in\{1,0,-1\}$ とする．
（$\pi_{10}$ が上の表の $\pi_{12}$，$\pi_{01}$ が $\pi_{21}$ に当たる．添字の付け方が年度で違うので注意．）
\[
\E[Z_i]=1\cdot\pi_{10}+(-1)\cdot\pi_{01}=\pi_{10}-\pi_{01},\qquad
\E[Z_i^2]=\pi_{10}+\pi_{01},\qquad
\V[Z_i]=(\pi_{10}+\pi_{01})-(\pi_{10}-\pi_{01})^2 .
\]
$\hat\pi_X-\hat\pi_Y=\bar Z$，$\sum_iZ_i=n_{10}-n_{01}$ なので $\V[\hat\pi_X-\hat\pi_Y]=\V[Z_i]/n$．
多項分布の $\V[n_{ij}]=n\pi_{ij}(1-\pi_{ij})$，$\Cov(n_{12},n_{21})=-n\pi_{12}\pi_{21}$ から直接計算しても同じ（\kakomon{2018}{2}〔1〕）：
\[
\V[\hat\delta]=\frac1{n^2}\{n\pi_{12}(1-\pi_{12})+n\pi_{21}(1-\pi_{21})+2n\pi_{12}\pi_{21}\}=\frac{\pi_{12}+\pi_{21}-(\pi_{12}-\pi_{21})^2}{n}.
\]

\subsection{検定統計量}
$H_0:\pi_{10}=\pi_{01}$（周辺の陽性率が等しい）の下では $\E[Z_i]=0$ なので $\V[Z_i]=\pi_{10}+\pi_{01}$．これを $(n_{10}+n_{01})/n$ で推定して
\[
W=\frac{(\hat\pi_X-\hat\pi_Y)^2}{\hat\V_0[\hat\pi_X-\hat\pi_Y]}=\frac{(n_{10}-n_{01})^2/n^2}{(n_{10}+n_{01})/n^2}=\frac{(n_{10}-n_{01})^2}{n_{10}+n_{01}} .
\]
中心極限定理と Slutsky の定理により，$H_0$ の下で $W$ は漸近的に $\chi^2_1$ に従う．
\paragraph{条件付き二項検定（\kakomon{2019}{3}〔2〕）}
不一致ペア数 $s$ を与えると，各不一致ペアが「処置1だけ陽性」である確率は $\pi_{12}/(\pi_{12}+\pi_{21})$ なので
$n_{12}\mid s\sim\Bin(s,\ \pi_{12}/(\pi_{12}+\pi_{21}))$，$H_0$ の下で $\Bin(s,1/2)$．
正規近似 $z=(n_{12}-s/2)/\sqrt{s/4}$ の2乗は $W$ に一致する．小標本では二項分布で正確な $P$ 値を求める．

\subsection{尤度比検定（\kakomon{2018}{2}〔3〕〔4〕）}
対数尤度 $\ell=\sum n_{ij}\log\pi_{ij}$ を $\sum\pi_{ij}=1$ の下で最大化すると $\hat\pi_{ij}=n_{ij}/n$．
$H_0:\pi_{12}=\pi_{21}$ の制約も加えたラグランジュ関数 $\ell-\lambda(\sum\pi_{ij}-1)-\phi(\pi_{12}-\pi_{21})$ を各 $\pi_{ij}$ で微分して0とおくと
\[
n_{11}=\lambda\pi_{11},\quad n_{12}=(\lambda+\phi)\pi_{12},\quad n_{21}=(\lambda-\phi)\pi_{21},\quad n_{22}=\lambda\pi_{22}.
\]
辺々足すと $n=\lambda\sum\pi_{ij}+\phi(\pi_{12}-\pi_{21})=\lambda$．$\pi_{12}=\pi_{21}$ から $n_{12}/(n+\phi)=n_{21}/(n-\phi)$，これを解いて
\[
\hat\pi^{(0)}_{11}=\frac{n_{11}}n,\qquad \hat\pi^{(0)}_{12}=\hat\pi^{(0)}_{21}=\frac{n_{12}+n_{21}}{2n},\qquad \hat\pi^{(0)}_{22}=\frac{n_{22}}n .
\]
一致セルの推定値は両モデルで同じなので尤度比では打ち消し合い，概要の $G^2$ を得る．
自由度は（対立仮説の自由パラメータ数3）$-$（帰無仮説の2）$=1$．

\subsection{条件付きオッズ比}
ペア $k$ ごとに異なるベースライン $\alpha_k$ をもつロジットモデル $\logit\Prob(Y_{k1}=1)=\alpha_k+\beta$，$\logit\Prob(Y_{k2}=1)=\alpha_k$ では，
一致ペアは $\beta$ の情報を持たず，不一致ペアが「処置1だけ陽性」である条件付き確率は $e^\beta/(1+e^\beta)$ で $\alpha_k$ が消える．
条件付き尤度は二項分布の形になり，$e^{\hat\beta}=n_{12}/n_{21}$．$H_0:\beta=0$ は条件付き二項検定そのもの．
マッチドペアの MH オッズ比（第7章）もこの値になる．

\begin{warnbox}
\begin{itemize}
\item ペアを独立な2群とみなして通常の $2\times2$ 表の $\chi^2$ 検定を行う（正の相関による分散の減少を無視し，通常は検出力を失う）．
\item 一致ペア $n_{11},n_{22}$ を統計量に入れる．分母を $n$ にする．
\item 2019年問3のように疾患群・非疾患群の表がある問題で，感度の比較なのに両方の表を使う（感度は疾患群の表だけ）．
\end{itemize}
\end{warnbox}

% ============================================================
\section{出題例}\label{sec:m06-shutsudai}
\begin{shutsudai}{\kakomon{2018}{2}：マッチドペアの有効率の比較（記号のみ，数値なし）}
〔1〕$\hat\delta=(n_{12}-n_{21})/n$ の分散．〔2〕正規近似による CI．
〔3〕$H_0:\delta=0$ と対立仮説の下の MLE（ラグランジュ乗数法）．〔4〕尤度比検定統計量と自由度1．
\end{shutsudai}
\begin{shutsudai}{\kakomon{2019}{3}〔2〕：2検査の感度の比較}
疾患群90人の2検査の表で，不一致は（A陽性・B陰性）7，（A陰性・B陽性）16の計23．
ヒントで $\Bin(23,0.5)$ が与えられ，$z=(7-11.5)/\sqrt{23/4}=-1.88$，$|z|<1.96$ で有意差なし．
\end{shutsudai}
\begin{shutsudai}{\kakomon{2024}{1}〔2〕〔3〕：$Z_i=X_i-Y_i$ から McNemar 統計量}
〔2-1〕$\E[Z_i]$，$\V[Z_i]$．〔2-2〕$H_0$ の下の分散推定を代入して $W=(n_{10}-n_{01})^2/(n_{10}+n_{01})$．
〔3〕200人で $n_{10}=26$，$n_{01}=16$：$W=100/42=2.38$，$\chi^2_1$ の5\%点3.84より小さく有意でない．
\end{shutsudai}

% ============================================================
\section{例題}\label{sec:m06-reidai}
\begin{reidai}[\kakomon{2024}{1}・\kakomon{2018}{2}型：対応のある2値データ]\label{reidai:m06-1}
ある疾患の患者250人に2種類の検査 $X$，$Y$ を行った結果を示す．
\begin{center}\small
\begin{tabular}{@{}lcc@{}}
\toprule
 & 検査 $Y$ 陽性 & 検査 $Y$ 陰性\\
\midrule
検査 $X$ 陽性 & 150 & 34\\
検査 $X$ 陰性 & 18 & 48\\
\bottomrule
\end{tabular}
\end{center}
$X_i,Y_i$ は陽性で1，陰性で0，$\pi_{jk}=\Prob(X_i=j,Y_i=k)$，$n_{jk}$ を $(X_i,Y_i)=(j,k)$ の人数，$Z_i=X_i-Y_i$ とする．
\begin{enumerate}[label=〔\arabic*〕]
\item $\E[Z_i]$，$\V[Z_i]$ を $\pi_{jk}$ で表せ．
\item $H_0:\pi_X=\pi_Y$（$\pi_X=\Prob(X_i=1)$，$\pi_Y=\Prob(Y_i=1)$）の検定統計量 $W=(\hat\pi_X-\hat\pi_Y)^2/\hat\V_0[\hat\pi_X-\hat\pi_Y]$ を $n_{jk}$ で表せ．
      分母は $H_0$ の下での分散の未知パラメータを推定値で置き換えたものとする．
\item $W$ が $H_0$ の下で近似的に従う分布を述べ，有意水準5\%で検定せよ．
\item 感度の差 $\delta=\pi_X-\pi_Y$ の推定値と，正規近似による95\%信頼区間を求めよ．
\item 尤度比検定統計量を求めよ．
\item 不一致ペアの数を与えた条件付き二項分布による $z$ 統計量を求め，$W$ との関係を述べよ．
\end{enumerate}
\end{reidai}

% ============================================================
\section{例題の詳しい解説}\label{sec:m06-kaisetsu}
\begin{kaitou}{reidai:m06-1}
$n_{11}=150$，$n_{10}=34$，$n_{01}=18$，$n_{00}=48$，$N=250$．

〔1〕$Z_i$ は $(X_i,Y_i)=(1,0)$ で1，$(0,1)$ で $-1$，一致なら0なので
\[
\E[Z_i]=\pi_{10}-\pi_{01},\qquad \E[Z_i^2]=\pi_{10}+\pi_{01},\qquad
\V[Z_i]=\pi_{10}+\pi_{01}-(\pi_{10}-\pi_{01})^2 .
\]

〔2〕$\hat\pi_X-\hat\pi_Y=\frac1N\sum Z_i=\frac{n_{10}-n_{01}}{N}$．
$\pi_X=\pi_{11}+\pi_{10}$，$\pi_Y=\pi_{11}+\pi_{01}$ なので $H_0$ は $\pi_{10}=\pi_{01}$ と同値で，その下では $\E[Z_i]=0$，$\V[Z_i]=\pi_{10}+\pi_{01}$．
これを $(n_{10}+n_{01})/N$ で推定すると $\hat\V_0[\hat\pi_X-\hat\pi_Y]=(n_{10}+n_{01})/N^2$ で
\[
W=\frac{(n_{10}-n_{01})^2/N^2}{(n_{10}+n_{01})/N^2}=\frac{(n_{10}-n_{01})^2}{n_{10}+n_{01}}\quad\text{（McNemar 統計量）}.
\]

〔3〕$H_0$ の下で $W$ は近似的に自由度1の $\chi^2$ 分布に従う．
$W=(34-18)^2/(34+18)=256/52=4.92>\chi^2_1(0.05)=3.841$ なので，有意水準5\%で $H_0$ を棄却する．
検査 $X$ の陽性率（$184/250=0.736$）は検査 $Y$（$168/250=0.672$）より高い．
患者だけの表なので，陽性率は感度であり，「$X$ のほうが感度が高い」と結論できる．

〔4〕$\hat\delta=(34-18)/250=0.064$．
\[
\hat\V[\hat\delta]=\frac{n_{10}+n_{01}}{N^2}-\frac{(n_{10}-n_{01})^2}{N^3}=\frac{52}{62500}-\frac{256}{15625000}=0.000832-0.0000164=0.0008156,
\]
$\SE=\sqrt{0.0008156}=0.02856$．95\%CI は $0.064\pm1.96\times0.02856=0.064\pm0.0560=(0.008,\,0.120)$．
区間が0を含まないので〔3〕と同じ結論．

〔5〕制約なしの MLE は $\hat\pi_{jk}=n_{jk}/N$，$H_0$ の下では一致セルは同じで不一致セルが $\hat\pi^{(0)}_{10}=\hat\pi^{(0)}_{01}=(34+18)/(2N)=26/N$．一致セルは打ち消し合い
\begin{align*}
G^2&=2\Bigl[34\log\frac{34}{26}+18\log\frac{18}{26}\Bigr]=2\Bigl[34\log\frac{68}{52}+18\log\frac{36}{52}\Bigr]\\
&=2[34\times0.2683+18\times(-0.3677)]=2(9.122-6.619)=5.01 .
\end{align*}
（付表より $\log(68/52)=\log1.3077=0.2683$，$\log(36/52)=\log0.6923=-0.3677$．）
自由度1の $\chi^2$ で $5.01>3.841$，有意．McNemar の $W=4.92$ とほぼ同じ値になる．

〔6〕不一致ペア $s=34+18=52$．$H_0$ の下で $n_{10}\mid s\sim\Bin(52,1/2)$，平均26，分散13．
\[
z=\frac{34-26}{\sqrt{13}}=\frac{8}{3.606}=2.22>1.96 .
\]
$z^2=64/13=4.92=W$ で，McNemar 検定は条件付き二項検定の正規近似（の2乗）である．
なお条件付きオッズ比は $34/18=1.89$．
\end{kaitou}

\begin{summarybox}
\begin{itemize}
\item 一致ペアは差に寄与しない．$\delta=\pi_{12}-\pi_{21}$．
\item $W=(n_{12}-n_{21})^2/(n_{12}+n_{21})\sim\chi^2_1$．$=z^2$，$z=(n_{12}-s/2)/\sqrt{s/4}$（条件付き $\Bin(s,1/2)$）．
\item $\V[\hat\delta]=\{\pi_{12}+\pi_{21}-(\pi_{12}-\pi_{21})^2\}/n$．LRT は不一致セルだけが残り自由度1．条件付き OR は $n_{12}/n_{21}$．
\end{itemize}
\end{summarybox}
